\section{Conclusion}
\label{sec:conclusion}

Geographic CVD completes an important practical choice in the CVD family:
BISECT can now split by graph-hop proximity or by projected-coordinate
proximity. The graph-distance mode is data-light and native to contiguity.
The geographic mode is useful when the map and the adjacency graph tell
different proximity stories.

\paragraph{When to use geographic CVD.}
Use \texttt{--cvd-metric geographic} when the split should be driven by map
distance: coastal states, peninsulas, long corridors, island-adjacent regions,
and sparse tracts with very different physical sizes. Use the default
\texttt{graph-distance} mode when centroid data is absent, when the graph
topology is the intended proximity model, or when the publication package has
not yet validated projection and centroid provenance.

\paragraph{Current command shape.}
\begin{verbatim}
algorithm:
  structure: centroidal-voronoi
  weights: geographic
  search: convergence
  cvd_iters: 20
  cvd_metric: geographic
  balance_tolerance: 0.5
\end{verbatim}

\paragraph{Legal and practitioner use.}
Geographic CVD should be presented as a reproducible compactness reference
plan, not as a legal compliance certificate. Compactness standards vary by
jurisdiction; Voting Rights Act analysis, communities of interest, incumbent
handling, county and municipal splits, and state-specific criteria remain
separate checks. The value of geographic CVD is that its proximity decision is
auditable from centroids, projection constants, base seed, and node path.

\paragraph{Remaining work.}
The next paper-quality and implementation milestones are clear:

\begin{enumerate}
\item Build archived NC/FL/WA benchmark packages with graph-distance and
  geographic CVD side by side.
\item Add ratio-aware CVD target populations so odd-seat nodes can split by
  $k_L:k_R$ rather than only by half population.
\item Record projection metadata in the run/audit manifest.
\item Validate the centroid fetch route and archive centroid file hashes.
\item Add state-specific projection dispatch only after tests and package
  metadata can prove which transform was used.
\end{enumerate}

\paragraph{Summary.}
T.10 explains CVD as a proximity-packing split primitive. T.11 makes the
geographic metric legible: show the coordinate field, show the ownership
frontier, show the snapped centroid update, and show the audit inputs needed
to reproduce the split. That is the standard this family needs before we move
from attractive geometry to durable public evidence.
